% Guide enseignant d'EP Manager (mod_ep et mod_epsynthesis).
% Compilation : latexmk -pdf guide-enseignant.tex
\documentclass[11pt,a4paper]{article}

% Version du guide : à mettre à jour à chaque modification, avec une ligne dans l'historique.
\newcommand{\titreguide}{Guide enseignant}
\newcommand{\versionguide}{1.1}
\newcommand{\dateguide}{novembre 2026}
\newcommand{\versionep}{1.0.4}
\newcommand{\versionsynthese}{1.0.0}

\usepackage{guide-commun}

\begin{document}

\pagedegarde{Responsable d'EP, enseignant référent : valider, noter, proposer un EP}{ens-synthese-a-traiter}

\tableofcontents

\section*{Historique des versions}
\addcontentsline{toc}{section}{Historique des versions}

\begin{center}
\begin{tabularx}{\linewidth}{@{}l l X@{}}
\toprule
\textbf{Version} & \textbf{Date} & \textbf{Modifications} \\
\midrule
1.0 & novembre 2026 & Première version, pour EP Manager 1.0. \\
1.1 & novembre 2026 & Séances du programme sans date (EP Manager 1.0.4). \\
\bottomrule
\end{tabularx}
\end{center}

\newpage

% ===========================================================================
\section{Vos rôles}

Un enseignant intervient dans l'enseignement personnalisé à trois titres, qui peuvent se
cumuler :
\begin{itemize}
  \item \textbf{responsable d'un EP du catalogue} : il inscrit et retire des étudiants, valide
    ou note les inscrits ;
  \item \textbf{enseignant référent} d'étudiants (le même que pour les stages) : il examine
    leurs déclarations d'EP hors catalogue et les éléments qu'ils demandent pour leur supplément
    au diplôme, à son étape du circuit fixé par la DEVE ;
  \item \textbf{auteur d'un EP} : il propose un nouvel EP ou met à jour le syllabus du sien.
\end{itemize}
Vous ne voyez que ce qui relève de vous : vos EP, vos étudiants référents.

\subsection{Une seule page pour toutes vos promotions}

L'activité \og Suivi de l'enseignement personnalisé\fg{} de votre cours d'équipe rassemble ce
que vous avez à faire dans toutes les promotions :
\begin{itemize}
  \item \chemin{Validation} : les demandes qui attendent votre décision, de la plus ancienne à
    la plus récente ;
  \item \chemin{Tableau de pilotage} : l'avancement de vos étudiants ;
  \item \chemin{À traiter} : les EP dont vous êtes responsable avec des inscrits à trancher,
    vos propositions et syllabus à reprendre, les éléments du supplément à examiner, que vous
    pouvez trancher directement sur la page.
\end{itemize}

\capture{ens-synthese-a-traiter}{Page \og À traiter\fg{} de la synthèse, toutes promotions
confondues.}

Chaque ligne renvoie vers la page de l'activité de la promotion concernée, décrite ci-dessous.

% ===========================================================================
\section{Responsable d'un EP du catalogue}

\subsection{Liste des participants}

Dans l'activité de la promotion, \chemin{Validation} commence par vos EP. Le bouton
\bouton{Participants} ouvre la liste des inscrits de l'un d'eux.

\capture{ens-participants}{Participants d'un EP : inscription en attente, note, décision.}

Depuis cette liste :
\begin{itemize}
  \item \textbf{valider ou noter} chaque inscrit : pour un EP noté, saisir la note (la
    validation suit le seuil fixé) ; sinon, choisir \og validé\fg{} ou \og non validé\fg{} ;
    ajouter un commentaire ; enregistrer toute la liste d'un coup ;
  \item \textbf{inscrire des étudiants} (\bouton{Inscrire des étudiants}) en cours de session, même après la fermeture des
    inscriptions ;
  \item \textbf{retirer} un étudiant, avec un motif ;
  \item \textbf{importer les notes} depuis un fichier CSV (la liste exportée de l'EP sert de
    modèle) : la page présente d'abord ce qui sera fait, ligne par ligne, puis applique sur
    confirmation.
\end{itemize}

\begin{astuce}
Si la DEVE a prévu une seconde étape (\og responsable puis DEVE\fg{}), votre validation fait
passer la demande à la DEVE ; sinon, elle est définitive et l'étudiant en est prévenu.
\end{astuce}

\subsection{Une demande à la fois}

Le bouton \bouton{Valider} d'une demande ouvre son détail : ce que l'étudiant a demandé, ses
justificatifs, l'historique des décisions. Vous y validez (en ajustant au besoin les ECTS
retenus), notez ou refusez avec un commentaire.

\capture{ens-validation}{Détail d'une demande, avec les boutons de décision.}

Si quelqu'un a tranché la même demande entre-temps, rien n'est enregistré et un message vous
prévient que la demande a changé.

% ===========================================================================
\section{Enseignant référent}

Les \textbf{déclarations} de vos étudiants référents (engagement, expérience
professionnelle, sport, académique externe\dots) arrivent dans \chemin{Validation}, rubrique
\og En attente de votre décision\fg{}. Vérifiez les justificatifs, puis validez (avec les ECTS
retenus) ou refusez en expliquant pourquoi.

Les \textbf{éléments du supplément au diplôme} demandés par vos étudiants arrivent dans
\chemin{Éléments de supplément à examiner} : catégorie, période, intitulé, preuves. Votre
validation transmet l'élément à la DEVE, qui décide en dernier ; votre commentaire lui est
transmis.

\capture{ens-supplement-items}{Éléments du supplément à examiner, avec les preuves.}

% ===========================================================================
\section{Proposer un EP, mettre à jour le sien}

\subsection{Proposer un nouvel EP}

\chemin{Propositions et syllabus} \textgreater{} \bouton{Proposer un nouvel EP} crée l'EP à l'état de
brouillon, dont vous devenez responsable. Renseignez :
\begin{itemize}
  \item ses caractéristiques : intitulé, type, ECTS, places, années concernées ;
  \item son \textbf{syllabus}, au modèle de la fiche ECTS : objectifs, acquis d'apprentissage,
    prérequis, contenu, méthodes d'enseignement, conditions d'examen, seconde session,
    bibliographie, langue, travail personnel ;
  \item son \textbf{programme} : les séances, avec leur nature (cours, TD, TP, clinique\dots),
    leur durée et leur intervenant (un enseignant du site, ou un intervenant extérieur par son
    nom). Les séances ne sont pas datées : le syllabus décrit l'EP, pas le calendrier d'une
    session. Le volume horaire total et par nature est calculé.
\end{itemize}

\capture{ens-syllabus-edit}{Rédaction du syllabus et du programme des séances.}

\capture{gu-ens-seances}{Programme des séances, sans date.}

Le bouton \bouton{Soumettre à la DEVE} transmet la proposition. La DEVE la publie (l'EP entre
au catalogue), vous la renvoie avec un commentaire (corrigez puis resoumettez) ou la refuse. Vous
êtes prévenu de sa décision.

\subsection{Mettre à jour un EP publié}

Depuis la même page, \bouton{Proposer une mise à jour} ouvre une \textbf{nouvelle version} du syllabus :
la version publiée reste au catalogue tant que la nouvelle n'est pas publiée. Selon le choix de
la DEVE, votre mise à jour est publiée directement ou passe par sa revue. Les versions
précédentes restent dans l'historique.

\capture{ens-propositions}{Propositions et syllabus : état de chaque version.}

La fiche de l'EP se télécharge en PDF depuis sa page \bouton{Syllabus}.

\end{document}
